\documentclass[runningheads]{llncs}
\usepackage[T1]{fontenc}
\usepackage[table]{xcolor}
\usepackage{graphicx}
\usepackage{booktabs}
\usepackage{array}
\usepackage{makecell}
\usepackage{amsmath,amssymb}
\usepackage{xurl}

\begin{document}

\title{Beyond Task Success: Measuring Workflow Fidelity in LLM-Based Agentic Payment Systems}

\author{
Donghao HUANG \inst{1,2}\orcidID{0009-0005-6767-4872}
\and
Joon Kiat CHUA \inst{1}\orcidID{0009-0001-8333-9792}
\and
Zhaoxia WANG \inst{1,\thanks{Corresponding author: zxwang@smu.edu.sg}}\orcidID{0000-0001-7674-5488}
}

%
\authorrunning{D. Huang, J.K. Chua, and Z. Wang}
% First names are abbreviated in the running head.
% If there are more than two authors, 'et al.' is used.
%

\institute{School of Computing and Information Systems, Singapore Management University, 80 Stamford Rd, Singapore 178902, Singapore \and Research and Development, Mastercard, 4250 Fairfax Dr, Arlington, VA 22203, USA}
\maketitle

\begin{abstract}
LLM-based multi-agent systems are increasingly deployed for payment workflows, yet prevailing metrics---Task Success Rate (TSR) and Agent Handoff F1-Score (HF1)---capture only final outcomes or unordered routing decisions. We introduce the \emph{Agentic Success Rate} (ASR), a trajectory-fidelity metric that compares observed and expected agent execution sequences at the \emph{transition} (bigram) level, decomposing performance into Transition Recall (completeness) and Transition Precision (efficiency). Applied to the Hierarchical Multi-Agent System for Payments (HMASP) across 18 LLMs and 90{,}000 task instances (5 repeats per model--scenario pair), ASR reveals that 11 of 18 models systematically skip a confirmation checkpoint during payment checkout---invisible to both TSR and HF1---while 7 models enforce the checkpoint perfectly. Notably, 12 open-weight models outperform both proprietary GPT models on average ASR. Prompt refinements and deterministic routing guards guided by ASR diagnostics yield substantial TSR improvements, with gains up to +93.8\,pp for previously struggling models, demonstrating that trajectory-level evaluation is essential in regulated domains.

\keywords{Agentic Success Rate \and Agentic Payments \and LLM Multi-Agent Systems \and Evaluation Metrics \and Digital Finance}
\end{abstract}

\section{Introduction}

The emergence of autonomous LLM-based agents has fueled rapid adoption of \textit{agentic payments}---using LLM agents to automate end-to-end payment workflows~\cite{mckinsey2025agenticpayments}. Major payment networks have launched agentic capabilities: Mastercard's Agent Pay~\cite{mastercard2025agentpay} and Visa's Intelligent Commerce~\cite{visa2025intelligentcommerce}, with projections reaching \$1.7 trillion by 2030~\cite{mckinsey2025agenticpayments}. In research, frameworks like HMASP~\cite{chua2025hmasp} have demonstrated the feasibility of hierarchical multi-agent systems for payments.

Despite these deployments, current evaluation frameworks remain inadequate. HMASP, for example, relies on Task Success Rate (TSR) and Agent Handoff F1-Score (HF1). However, TSR captures only final outcomes, and HF1 evaluates routing in a bag-of-edges fashion, ignoring temporal order. Neither captures \textit{workflow fidelity}---the agent's adherence to the expected sequence of steps. In payment systems, auditability is critical~\cite{pcisecurity2025ai}, and the sequential path to the result is as important as the outcome itself~\cite{chen2025benchmarksfail,cemri2025multifail}.

Drawing on process mining conformance checking~\cite{vanderaalst2012conformance,carmona2018conformance}, we propose the \textbf{Agentic Success Rate (ASR)}, a trajectory-level metric that compares agent-to-agent \emph{transitions} (consecutive handoff pairs) between observed and expected workflows.
Our contributions are: (1)~ASR as a multi-agent payment metric; (2)~evaluation of 18 LLMs on 90{,}000 instances; (3)~discovery of workflow shortcuts invisible to TSR/HF1; and (4)~demonstration that ASR-guided engineering yields dramatic TSR improvements for previously struggling models.

\section{Related Work}

\noindent\textbf{LLM-Based MAS for Payments.}
Multi-agent systems have been applied to software engineering~\cite{hong2024metagpt}, trading~\cite{tradingagents2024}, and investment~\cite{yu2024fincon}. In payments, prior work addressed process automation~\cite{guan2024llmprocessautomation} and fraud detection~\cite{dahiphale2024llmscampayments} but not end-to-end agentic workflows. HMASP~\cite{chua2025hmasp} filled this gap with a hierarchical agent framework.

\noindent\textbf{Evaluating Agentic Systems.}
Standard evaluation focuses on task completion~\cite{liu2023agentbench,kapoor2024ai}, but binary metrics obscure behavioral nuances. AgentBoard~\cite{agentboard2024} introduced Progress Rate; WorFEval~\cite{worfeval2025} uses subsequence matching for workflow planning; $\tau$-bench~\cite{taubench2024} proposed pass$^k$ for reliability; and Chen et al.~\cite{chen2025benchmarksfail} argued standard benchmarks underestimate financial risks. ASR builds on this line of work~\cite{mohammadi2025agenteval}, grounding trajectory evaluation in process mining conformance checking~\cite{vanderaalst2012conformance,carmona2018conformance}.% Specifically, it adapts bigram overlap---used in software process analysis~\cite{cook1999softwareprocess} and NLP evaluation~\cite{lin2004rouge}---to multi-agent payment workflows.

\section{Methodology}

% \subsection{Background: HMASP Architecture}

% HMASP~\cite{chua2025hmasp} organizes agents into four levels: (1)~the CPA as entry point; (2)~Supervisor agents for domain workflows; (3)~Routing agents for task initiation; and (4)~Process Summary agents for outcome responses. It supports card registration (T1), card retrieval (T2), payment processing (T3), and irrelevant input rejection (T4).

\subsection{Agentic Success Rate (ASR)}

ASR compares the full ordered trajectory against the expected one at the \emph{transition} level, drawing on the conformance checking paradigm from process mining~\cite{vanderaalst2012conformance,carmona2018conformance}. Each consecutive pair of agents in a trajectory is a transition (bigram); for example, the trajectory $[A, B, C]$ yields the transition multiset $\{(A{,}B),\;(B{,}C)\}$. Given expected trajectory $E = [e_1, \ldots, e_n]$ and observed trajectory $O = [o_1, \ldots, o_m]$, let $B_E$ and $B_O$ denote their respective transition multisets:
%
\begin{align}
\text{TR} &= \frac{|B_E \cap B_O|}{|B_E|}, \quad
\text{TP} = \frac{|B_E \cap B_O|}{|B_O|}, \quad
\text{ASR} = F_1(\text{TR}, \text{TP})
\end{align}
%
An ASR of 100\% indicates perfect transition alignment. Low Transition Recall (TR) signals skipped handoffs (``model moves'' in process mining terminology); low Transition Precision (TP) signals redundant ones (``log moves''). Because ASR operates on transitions rather than individual agent visits or unordered edge sets, it directly measures whether each agent handed off control to the correct next agent---precisely the property needed for compliance auditing---and generalizes to arbitrary agentic architectures.

\subsection{Evaluation Setup}

% We evaluate 18 models (11 from HMASP~\cite{chua2025hmasp} plus 7 new) spanning open-weight and proprietary LLMs on 1{,}000 data points across 4 scenarios (250 each), yielding 18{,}000 instances. We report TSR, HF1, and ASR. System engineering comprises iterative prompt refinement and two deterministic routing guards addressing failure modes that prompting alone cannot resolve.

\noindent\textbf{Models.}
We evaluated 18 models spanning both open-weight and proprietary LLMs. From the original HMASP study~\cite{chua2025hmasp}, we include Qwen3 (8B, 14B, 32B), Qwen2.5 (7B, 14B, 32B), Mistral-Small-3.2 (24B), Magistral (24B), Llama~3.1 (8B, 70B), and GPT-4.1. Additionally, we evaluated seven models not included in the original study: GPT-5.2, GPT-OSS (20B, 120B), and Gemma4 (E2B, E4B, 26B, 31B). Proprietary models (GPT-4.1, GPT-5.2) are accessed via the official OpenAI API, while all open-weight models are run locally using Ollama v0.20.0 on an Apple M5~Max with 128\,GB unified memory and 4-bit quantization (MXFP4 for GPT-OSS models; Q4\_K\_M for all others).

\noindent\textbf{Dataset and Repeat Protocol.}
We use the HMASP evaluation dataset of 1{,}000 data points across four scenarios (250 each): \textbf{T1}~(card registration; e.g., \textit{``Add a new card''}), \textbf{T2}~(card retrieval; e.g., \textit{``Show my saved cards''}), \textbf{T3}~(payment processing; e.g., \textit{``Complete my purchase''}), and \textbf{T4}~(irrelevant input rejection; e.g., \textit{``What's the weather today?''}). Each model--scenario pair is evaluated 5 times independently, yielding 90{,}000 total instances. We report the mean and standard deviation across the 5 repeats.

\noindent\textbf{Metrics.}
We report three metrics, the first two from the original HMASP evaluation~\cite{chua2025hmasp}: \textbf{(1)~Task Success Rate (TSR):} the proportion of runs completing the correct workflow with correct information saved or retrieved; \textbf{(2)~Agent Handoff F1-Score (HF1):} the F1-score comparing actual vs.\ expected agent-to-agent handoffs; and \textbf{(3)~Agentic Success Rate (ASR):} trajectory-level workflow fidelity as defined above.

\noindent\textbf{System Engineering.}
We improve the HMASP system through two complementary techniques, both guided by ASR diagnostics. \textbf{(i)~Prompt engineering:} agent prompts are refined iteratively to better specify roles, expected handoffs, and task boundaries. \textbf{(ii)~Deterministic routing guards:} a \textit{cart-items guard} in the payment supervisor suppresses the payment-review path when \texttt{cart\_items} is empty, and \textit{card-view pattern matching} in the CPA provides keyword-based fallback routing for ambiguous card intents.

\section{Results and Discussion}

\subsection{Effectiveness of System Engineering}

Table~\ref{tab:paper_vs_latest} compares TSR between HMASP~\cite{chua2025hmasp} and our system for the three most-improved models (ranked by average $\Delta$ across all tasks). The gains are striking: Llama3.1:8b improves by an average of +67.9\,pp, Magistral:24b by +54.2\,pp, and Llama3.1:70b by +33.5\,pp---confirming that failures were addressable through engineering rather than model replacement. %Top HMASP baselines (GPT-4.1, Qwen2.5:32b) were already near-saturated; notably, Qwen2.5:32b now achieves perfect TSR (100\%) with zero variance across all four task types.

\begin{table}[!ht]
    \centering
    \caption{HMASP vs.\ Ours: TSR (\%) for the three most-improved models. H~=~HMASP, O~=~Ours (mean over 5 repeats), $\Delta$~= O - H. Bold: $\geq$20\,pp gains.}
    \label{tab:paper_vs_latest}
    \renewcommand{\arraystretch}{0.85}
    \setlength{\tabcolsep}{1.5mm}
    \scriptsize
    \begin{tabular}{l rr@{\;}r rr@{\;}r rr@{\;}r rr@{\;}r}
        \toprule
        & \multicolumn{3}{c}{\textbf{T1}} & \multicolumn{3}{c}{\textbf{T2}} & \multicolumn{3}{c}{\textbf{T3}} & \multicolumn{3}{c}{\textbf{T4}} \\
        \cmidrule(lr){2-4} \cmidrule(lr){5-7} \cmidrule(lr){8-10} \cmidrule(lr){11-13}
        \textbf{Model} & \textbf{H} & \textbf{O} & $\boldsymbol{\Delta}$ & \textbf{H} & \textbf{O} & $\boldsymbol{\Delta}$ & \textbf{H} & \textbf{O} & $\boldsymbol{\Delta}$ & \textbf{H} & \textbf{O} & $\boldsymbol{\Delta}$ \\
        \midrule
        Llama3.1:8b & 6.0 & 99.8 & \textbf{\scriptsize+93.8} & 10.4 & 98.8 & \textbf{\scriptsize+88.4} & 4.0 & 93.4 & \textbf{\scriptsize+89.4} & 100 & 100 & \scriptsize$+$0.0 \\
        Magistral:24b & 44.0 & 90.0 & \textbf{\scriptsize+46.0} & 18.4 & 97.7 & \textbf{\scriptsize+79.3} & 3.2 & 94.7 & \textbf{\scriptsize+91.5} & 100 & 100 & \scriptsize$+$0.0 \\
        Llama3.1:70b & 66.0 & 100 & \textbf{\scriptsize+34.0} & 74.0 & 100 & \textbf{\scriptsize+26.0} & 26.0 & 98.9 & \textbf{\scriptsize+72.9} & 98.8 & 100 & \scriptsize+1.2 \\
        \bottomrule
    \end{tabular}
\end{table}

\subsection{ASR Reveals Hidden Workflow Deviations}

Table~\ref{tab:combined} reports T3 metrics (mean\,$\pm$\,std over 5 repeats) and average across T1--T4 for all 18 models, sorted by average ASR. Seven models achieve a perfect 100\% on all three T3 metrics: all four Gemma4 variants, both GPT-OSS models, and \texttt{MSmall3.2:24b}. A striking result is that \textbf{12 open-weight models outperform the proprietary GPT-4.1 and GPT-5.2} on average ASR. Among the 11 models that exhibit a T3 ASR gap, the range spans from 0.01\% (\texttt{Qwen2.5:32b}) to 5.45\% (\texttt{Qwen2.5:7b}); notably, \texttt{GPT-4.1} and \texttt{GPT-5.2} show the largest gaps among high-performing models at 2.90\% and 2.72\% respectively. The \colorbox{yellow!30}{highlighted} cells are the sharpest examples: \texttt{Qwen2.5:32b}, \texttt{Qwen3:32b}, and \texttt{Qwen3:8b} achieve T3 TSR\,=\,HF1\,=\,100\% yet ASR\,$<$\,100\%, exposing skipped confirmation checkpoints invisible to conventional metrics.

\begin{table}[!ht]
    \centering
    \caption{T3 (Payment) results (\%, mean\,$\pm$\,std over 5 repeats) and average across T1--T4 for all 18 models. Sorted by Avg ASR desc. \colorbox{yellow!30}{Yellow}: TSR\,=\,HF1\,=\,100\% but ASR\,$<$\,100\%.}
    \label{tab:combined}
    \renewcommand{\arraystretch}{0.9}
    \setlength{\tabcolsep}{1mm}
    \scriptsize
    \begin{tabular}{l ccc @{\hskip 3mm} c@{\hskip 1.5mm}c@{\hskip 1.5mm}c}
        \toprule
        & \multicolumn{3}{c}{\textbf{T3 (Payment)}} & \multicolumn{3}{c}{\textbf{Avg (T1--T4)}} \\
        \cmidrule(lr){2-4} \cmidrule(lr){5-7}
        \textbf{Model} & \textbf{TSR} & \textbf{HF1} & \textbf{ASR} & \textbf{TSR} & \textbf{HF1} & \textbf{ASR} \\
        \midrule
        \textit{6 models}$^\dagger$ & 100 & 100 & 100 & 100.0 & 100.0 & 100.0 \\
        Qwen2.5:32b      & 100 & 100 & \cellcolor{yellow!30}99.99$\pm$0.02 & 100.0 & 100.0 & 100.0 \\
        Qwen3:32b        & 100 & 100 & \cellcolor{yellow!30}99.97$\pm$0.04 & 100.0 & 100.0 & 100.0 \\
        Qwen3:14b        & 99.7$\pm$0.3  & 99.8$\pm$0.2  & 99.6$\pm$0.3  & 99.9  & 100.0 & 99.9 \\
        Qwen3:8b         & 100 & 100 & \cellcolor{yellow!30}99.97$\pm$0.02 & 99.9  & 100.0 & 99.9 \\
        Llama3.1:70b     & 98.9$\pm$1.2  & 98.8$\pm$1.1  & 98.8$\pm$1.2  & 99.7  & 99.7  & 99.7 \\
        Qwen2.5:14b      & 99.8$\pm$0.2  & 99.8$\pm$0.2  & 99.6$\pm$0.2  & 99.3  & 99.6  & 99.6 \\
        GPT-4.1          & 99.8$\pm$0.4  & 99.9$\pm$0.2  & 97.1$\pm$0.3  & 99.5  & 99.7  & 98.8 \\
        GPT-5.2          & 99.9$\pm$0.2  & 100.0$\pm$0.1 & 97.3$\pm$0.1  & 99.4  & 99.7  & 98.8 \\
        Gemma4:e2b       & 100 & 100 & 100 & 98.4  & 98.8  & 98.4 \\
        Llama3.1:8b      & 93.4$\pm$0.9  & 96.1$\pm$0.7  & 93.3$\pm$1.0  & 98.0  & 98.8  & 98.0 \\
        Magistral:24b    & 94.7$\pm$1.2  & 97.1$\pm$0.6  & 93.9$\pm$1.0  & 95.6  & 97.7  & 95.5 \\
        Qwen2.5:7b       & 53.3$\pm$2.1  & 64.2$\pm$2.3  & 47.8$\pm$2.0  & 77.5  & 85.1  & 77.1 \\
        \bottomrule
        \multicolumn{7}{l}
        {\scriptsize $^\dagger$Gemma4:e4b/26b/31b, GPT-OSS:20b/120b, MSmall3.2:24b.} \\
    \end{tabular}
\end{table}

The root cause is consistent across all 11 models exhibiting a T3 ASR gap. The expected T3 trajectory has 11 agent hops (10 transitions), including an intermediate return to the CPA for explicit user \emph{confirmation} before payment processing. For direct payment phrases (e.g., ``Process my payment''), models infer unambiguous intent and skip this checkpoint, executing 9 hops (8 transitions). Of 10 expected transitions, 8 appear: TR\,$= 8/10 = 80\%$, TP\,$= 8/8 = 100\%$, ASR\,$= F_1(80\%, 100\%) = 88.9\%$ per affected sample. Crucially, every single hidden deviation across all affected models follows this \textit{identical} 9-step shortcut---there is exactly one deviant trajectory pattern, reflecting a systematic interaction between input phrasing and model reasoning rather than stochastic variation.
% This failure mode reflects \textit{missing transitions} (low TR) rather than redundant ones. Payments still succeed because the skipped steps serve orchestration and compliance rather than functional transaction requirements; TSR and HF1 miss the deviation because the correct handoff types all appear, just not in the expected order. 
This has regulatory significance: PCI-DSS~\cite{pcisecurity2025ai} requires auditable workflows, and bypassing confirmation checkpoints creates gaps in the audit trail even when outcomes are correct.% The seven zero-deviation models prove that the CPA confirmation checkpoint can be consistently enforced without sacrificing task performance.

\section{Conclusion}

We introduced ASR as a trajectory-level evaluation metric for LLM-based multi-agent payment systems. Across 18 LLMs and 90{,}000 instances, ASR exposes a systematic procedural deviation: 11 of 18 models skip a CPA-mediated confirmation checkpoint during payment checkout (T3 ASR gaps of 0.01--6.7\%) despite near-perfect TSR and HF1, while 7 models enforce the checkpoint perfectly. Notably, 12 open-weight models outperform both proprietary GPT models on average ASR, with \texttt{GPT-4.1} and \texttt{GPT-5.2} showing the largest gaps among high-performing models. ASR-guided system engineering yields dramatic improvements for previously struggling models (up to +93.8\,pp). As payment networks deploy agentic capabilities at scale~\cite{mastercard2025agentpay,visa2025intelligentcommerce}, ASR provides the visibility needed to ensure procedural fidelity---essential where compliance, auditability, and liability attribution are at stake. Future work includes validating ASR in production systems, weighted variants for security-critical transitions, architectural enforcement of the T3 checkpoint, and cross-architecture evaluations~\cite{visa2025intelligentcommerce}.

\bibliographystyle{splncs04}
\bibliography{references}

\end{document}
